\chapter{Requirements Analysis and Specification}

\section{Identification of actors}

The EAM platform distinguishes four main roles:
\begin{itemize}
\item \textbf{Administrator (Admin)}: manages users, machines, documents in the knowledge base, and has an overall view of the platform;
\item \textbf{Operations Manager (ChefOp)}: consultation and overview of labels and statuses;
\item \textbf{Head Technician (ChefTech)}: plans interventions, oversees work orders, arbitrates priorities, and reviews reliability dashboards;
\item \textbf{Technician (Thech)}: executes intervention orders, views the health status of the machines assigned to them, and interacts with the conversational assistant to get help in the field.
\end{itemize}

\noindent
Figure~\ref{fig:use-case-diagram} presents the detailed use case diagram, showing the interactions between the four roles and the complete set of system features.

\begin{figure}[H]
\centering
\begin{tikzpicture}[scale=0.74,
  uc/.style={ellipse, draw, align=center, font=\footnotesize,
             text width=2.8cm, inner sep=3pt},
  uca/.style={ellipse, draw, align=center, font=\scriptsize,
              text width=2.0cm, inner sep=2pt},
  inc/.style={dashed, -{Latex}},
  ext/.style={dashed, -{Latex}},
]

%% ---- SYSTEM BOUNDARY ----
\draw[thick] (3.5,-2.2) rectangle (19.0,22.5);
\node[above, font=\small\bfseries] at (11.25,22.5) {EAM APP};

%% ---- ACTOR: Chefop (centered at y≈20.8) ----
\draw (1.5,21.4) circle (0.32);
\draw (1.5,21.08) -- (1.5,20.35);
\draw (1.1,20.72) -- (1.9,20.72);
\draw (1.5,20.35) -- (1.2,19.75) (1.5,20.35) -- (1.8,19.75);
\node[below, font=\footnotesize] at (1.5,19.75) {Chefop};

%% ---- ACTOR: Thech (centered at y≈13.5) ----
\draw (1.5,14.1) circle (0.32);
\draw (1.5,13.78) -- (1.5,13.05);
\draw (1.1,13.42) -- (1.9,13.42);
\draw (1.5,13.05) -- (1.2,12.45) (1.5,13.05) -- (1.8,12.45);
\node[below, font=\footnotesize] at (1.5,12.45) {Thech};

%% ---- ACTOR: ChefTech (centered at y≈6.5) ----
\draw (1.5,7.1) circle (0.32);
\draw (1.5,6.78) -- (1.5,6.05);
\draw (1.1,6.42) -- (1.9,6.42);
\draw (1.5,6.05) -- (1.2,5.45) (1.5,6.05) -- (1.8,5.45);
\node[below, font=\footnotesize] at (1.5,5.45) {ChefTech};

%% ---- ACTOR: Admin (centered at y≈0.5) ----
\draw (1.5,1.1) circle (0.32);
\draw (1.5,0.78) -- (1.5,0.05);
\draw (1.1,0.42) -- (1.9,0.42);
\draw (1.5,0.05) -- (1.2,-0.55) (1.5,0.05) -- (1.8,-0.55);
\node[below, font=\footnotesize] at (1.5,-0.55) {Admin};

%% ---- USE CASES (main column x=8.5) ----
\node[uc] (eth)    at (8.5, 21.5) {manage\\labels};
\node[uc] (stat)   at (8.5, 19.6) {view machine\\status};
\node[uc] (bt)     at (8.5, 17.6) {manage\\intervention orders};
\node[uc] (plan)   at (8.5, 15.6) {view\\schedule};
  \node[uca] (jr)   at (5.5, 13.4) {by day};
  \node[uca] (sem)  at (8.5, 13.4) {by week};
  \node[uca] (mois) at (11.5,13.4) {by month};
\node[uc] (rap)    at (8.5, 11.5) {view\\reports};
\node[uc] (arch)   at (8.5,  9.7) {manage archive};
\node[uc] (ot)     at (8.5,  7.8) {manage\\work order tasks};
\node[uc] (alerte) at (8.5,  5.9) {manage\\alerts};
  \node[uca] (prio) at (5.5,  4.0) {by priority};
\node[uc] (gplan)  at (8.5,  2.3) {manage\\schedule};
\node[uc] (roles)  at (8.5,  0.5) {manage\\roles/users};

%% ---- AUTHENTICATION (right side, large ellipse) ----
\node[uc, text width=3.2cm, minimum height=3.2cm]
  (auth) at (16.0, 11.0) {authentication};

%% ---- ASSOCIATIONS actor -> use case ----
\draw (1.5,21.4) -- (eth);
\draw (1.5,21.4) -- (stat);
\draw (1.5,14.1) -- (bt);
\draw (1.5,14.1) -- (plan);
\draw (1.5,14.1) -- (rap);
\draw (1.5,14.1) -- (arch);
\draw (1.5,7.1)  -- (ot);
\draw (1.5,7.1)  -- (alerte);
\draw (1.5,1.1)  -- (gplan);
\draw (1.5,1.1)  -- (roles);

%% ---- <<extends>> RELATIONS ----
\draw[ext] (plan) -- node[midway, above left,  font=\tiny]{<<extends>>} (jr);
\draw[ext] (plan) -- node[midway, right,        font=\tiny]{<<extends>>} (sem);
\draw[ext] (plan) -- node[midway, above right,  font=\tiny]{<<extends>>} (mois);
\draw[ext] (alerte) -- node[midway, left, font=\tiny]{<<extends>>} (prio);

%% ---- <<includes>> RELATIONS to authentication ----
%% A single label on the middle arrow for readability
\draw[inc] (eth)    -- (auth);
\draw[inc] (stat)   -- (auth);
\draw[inc] (bt)     -- (auth);
\draw[inc] (plan)   -- node[midway, above, sloped, font=\tiny]{<<includes>>} (auth);
\draw[inc] (rap)    -- (auth);
\draw[inc] (arch)   -- (auth);
\draw[inc] (ot)     -- (auth);
\draw[inc] (alerte) -- (auth);
\draw[inc] (gplan)  -- (auth);
\draw[inc] (roles)  -- (auth);

\end{tikzpicture}
\caption{Use case diagram of the EAM platform}
\label{fig:use-case-diagram}
\end{figure}

\section{Functional requirements}

\subsection{Authentication and access management}

The system must require authentication with a username and password before granting any access to the platform.
Each user is assigned a single role (Administrator, Operations Manager, Head Technician, Technician) that determines the full set of accessible features and data.
Any attempt to access an unauthorized resource must be rejected by the server, independently of any interface-side control.

\subsection{User and role management}

The administrator must be able to create, edit, deactivate, and delete user accounts, and assign them a role.
The system must guarantee that a deactivated account can no longer authenticate, without deleting its associated action history.

\subsection{Machine asset management}

The system must allow machines to be registered with their technical characteristics (identifier, category, location, commissioning date).
Each machine has a detailed, viewable record combining its current health status, maintenance history, active alerts, and ongoing predictions.
The administrator can add, edit, or archive a machine without affecting the existing work orders attached to it.

\subsection{Machine labeling}

The Operations Manager must be able to attach descriptive labels to each machine to facilitate classification and search.
The system must support creating, editing, and deleting labels, as well as assigning them to several machines at once.

\subsection{Work order management}

The system must support creating a work order (OT) by specifying the machine concerned, the nature of the intervention, the priority, and the desired date.
A work order follows a defined lifecycle: pending, assigned, in progress, completed, or cancelled. Every transition must be logged with the identity of the user and a timestamp.
The Head Technician can view all work orders, filter them by status or priority, and reassign them in the event a technician is unavailable.

\subsection{Intervention order management}

An intervention order (BT) is created from a work order and assigned to a specific technician with a time slot.
The technician must be able to view their assigned intervention orders, record the actions performed and the parts replaced, and close the intervention.
When an intervention is closed, the system must automatically trigger an update of the machine's health score and a retroactive evaluation of earlier predictions.

\subsection{Intervention scheduling}

The system must provide a calendar view of planned interventions, navigable by day, week, and month.
The Head Technician must be able to schedule a new intervention, modify an existing time slot, and check technician availability before any assignment.
The schedule must visually flag assignment conflicts (a technician already booked for the same time slot).

\subsection{Telemetry data collection and monitoring}

The system must record sensor readings for each machine: air temperature, process temperature, rotational speed, torque, and tool wear.
This data constitutes the primary source for the prediction models and must be viewable as a timestamped history per machine.
In the absence of real telemetry data for a machine, the system must explicitly report the unavailability of predictions rather than substituting default values.

\subsection{Machine learning predictions}

For each machine with telemetry data available, the system must expose six predictive indicators computed by machine learning models:
\begin{itemize}
  \item \textbf{Failure probability (P1)}: estimate of short-term failure risk;
  \item \textbf{Likely failure type (P2)}: classification among known failure modes;
  \item \textbf{Remaining useful life — RUL (P3)}: estimate in days until the next required intervention, together with a confidence interval;
  \item \textbf{Behavioral anomaly score (P4)}: detection of abnormal sensor behavior via an ensemble of four methods (Isolation Forest, Z-score, cluster deviation, autoencoder);
  \item \textbf{Processing priority (P5)}: recommended priority level for the associated work order;
  \item \textbf{Maintenance schedule (P6)}: proposed optimal date for the next preventive intervention.
\end{itemize}
Each prediction must be accompanied by a plain-language explanation, free of technical jargon, accessible to a non-specialist user.

\subsection{Alerts and notifications}

The system must automatically generate alerts when a model detects a high failure risk, a behavioral anomaly, or a foreseeable parts shortage.
Alerts must be classified by criticality level and viewable by the Head Technician, with filtering by priority.
An alert that has already been handled (linked to a closed work order) must be automatically archived and no longer appear in the active list.

\subsection{Spare parts coordination}

Based on failure predictions, the system must estimate future spare parts needs per machine and generate provisioning proposals.
These proposals take the form of draft purchase orders subject to explicit human validation: no order can be issued automatically without approval from the Head Technician.
The system also exposes an operational readiness score per machine, combining health status, stock level, and proximity to the next intervention.

\subsection{Knowledge base and conversational assistant}

The administrator must be able to import maintenance documents (manuals, technical data sheets, procedures) into the platform's knowledge base.
The system must allow any authenticated user to query this knowledge base in natural language via a conversational assistant.
The assistant's responses must combine information extracted from relevant documents with the real-time health status of the machine in question, in order to provide contextualized advice to the technician in the field.

\subsection{Dashboards and reports}

Each role has a personalized dashboard: a fleet-wide overview for the Administrator and the Operations Manager, an intervention and alert monitoring view for the Head Technician, and an assigned-tasks view for the Technician.
The system must support viewing summary reports on maintenance activity, machine performance, and the evolution of predictive indicators.

\subsection{Archiving and traceability}

Every action performed on a work order, an intervention order, or an alert must be logged in an audit trail with the identity of the author, the date, and the nature of the action.
Closed interventions and handled alerts must be preserved in a viewable, non-modifiable history, serving as the basis for the continuous evaluation of the predictive models.

\section{Non-functional requirements}

\subsection{Application security}

Access to the platform must be protected by token-based authentication (JWT), renewed at each session.
Authorization control must be enforced server-side on every API endpoint, independently of the user interface: no sensitive data must be accessible through simple manipulation of the URL or the HTTP request.
Passwords must be stored in hashed form; application secrets (API keys, connection parameters) must be isolated in environment variables and must never appear in the source code.
Every significant action (creating, editing, or closing a work order or intervention order, validating an alert) must be logged in a non-modifiable audit trail.

\subsection{DevSecOps pipeline}

Security must be built into the development process through a dedicated GitLab CI/CD pipeline, covering the entire software supply chain.
The pipeline includes the following stages, run automatically on every code push:

\begin{itemize}
  \item \textbf{Secret scanning (Gitleaks)}: detection of any secret or sensitive credential accidentally present in the source code or the Git history;
  \item \textbf{Software composition analysis (SCA)}: audit of Python (\textit{pip-audit}) and JavaScript (\textit{pnpm audit}) dependencies to identify known vulnerabilities (CVEs) and non-compliant licenses;
  \item \textbf{Static application security testing (SAST)}: source code inspection by SonarQube Community Edition (self-hosted instance), covering injection vulnerabilities, error handling, cryptographic quality, and data exposure;
  \item \textbf{Docker image and infrastructure analysis (Trivy)}: scanning of the five Docker images produced by the build, as well as the Docker Compose files, looking for CVEs in system layers and infrastructure misconfigurations;
  \item \textbf{Dynamic application security testing (DAST — OWASP ZAP)}: active scanning of the running application via the ZAP automation framework, covering authenticated contexts (four roles: Administrator, Head Technician, Operations Manager, Technician) and unauthenticated services (ML microservice, RAG service, frontend), targeting OWASP Top 10 vulnerabilities such as SQL injection, XSS, HTTP header misconfiguration, and authentication flaws.
\end{itemize}

The results of each scanner are consolidated by a centralized reporting script (\texttt{security\_report.py}) that produces a structured terminal report on every pipeline run.
All security stages are configured in non-blocking mode (\textit{allow\_failure: true}), providing immediate visibility into results without blocking delivery, while keeping the findings in a report available for review after every run.

\subsection{Performance}

The response time for common operations (viewing a machine record, listing work orders, updating a status) must remain under two seconds under normal usage conditions.
Machine learning model inference, which involves several computation pipelines, must complete in under five seconds for all six indicators of a machine.
The interface must remain responsive during asynchronous calls to the ML microservice, displaying a loading indicator rather than blocking the user.

\subsection{Availability and resilience}

Classic management features — work orders, intervention orders, scheduling — must remain fully operational even if the Machine Learning microservice is unavailable.
If telemetry data is missing for a machine, the system must explicitly report it and provide a health score computed from maintenance history, without substituting fictitious values.
The platform's various services (backend, ML microservice, RAG service, object storage) must be deployed as independent containers, allowing any one of them to be restarted or updated without interrupting the others.

\subsection{Maintainability and extensibility}

The architecture must clearly separate responsibilities between the presentation, business logic, and data access layers, to facilitate adding new predictive models or new roles without a full overhaul.
Database migrations must be versioned and applied incrementally through a dedicated tool (Alembic), ensuring traceability of schema evolution.
The ML model code must be organized so that an individual model can be replaced or updated without affecting the other prediction pipelines.

\subsection{Portability and deployment}

The entire platform must be able to run on any environment with Docker and Docker Compose, without any manual dependency installation.
Environment-specific configuration (URLs, keys, connection parameters) must be centralized in a single \texttt{.env} file, kept separate from the source code.

\subsection{Ergonomics and usability}

The interface must adapt its menus, dashboards, and available actions to the role of the logged-in user, without exposing features beyond their scope.
Information derived from the predictive models must be presented in business language — failure probability, remaining useful life, parts needed — without algorithmic terminology (model names, raw statistical metrics).
The interface must be usable from a desktop workstation as well as from a tablet used in the field, with a layout adapted to each screen format.

\subsection{Data integrity and confidentiality}

Data for each machine must be accessible only to users whose role explicitly authorizes it; a technician must not be able to view data for a machine that is not assigned to them.
The history of interventions, alerts, and predictions must be preserved with integrity: no retroactive modification is allowed once an order or an alert has been closed.

\section{Use case diagram}

\begin{table}[h]
\centering
\begin{tabularx}{\textwidth}{lX}
\toprule
\textbf{Actor} & \textbf{Main use cases} \\
\midrule
Administrator & Manage users, manage machines, manage the knowledge base, view global statistics \\
Head Technician & Plan interventions, oversee work orders, arbitrate priorities, review reliability dashboards, validate provisioning proposals \\
Technician & View assigned intervention orders, close an intervention, view a machine's health status, query the conversational assistant \\
\bottomrule
\end{tabularx}
\caption{Use cases by actor}
\label{fig:cas-utilisation}
\end{table}
